\section{Background}
\label{sec:background}

\subsection{Differential privacy and Gaussian mechanisms}

We use the standard definition of approximate differential
privacy~\cite{dwork2006calibrating}: a randomised mechanism
$\mathcal{M}\colon \mathcal{Z}^* \to \mathcal{O}$ satisfies $(\eps,
\delta)$-differential privacy ($(\eps, \delta)$-DP) if, for any pair of
adjacent databases $D \sim D'$ (differing on a single record) and every
measurable $S \subseteq \mathcal{O}$,
\[
\Pr[\mathcal{M}(D) \in S] \;\leq\; e^\eps \Pr[\mathcal{M}(D') \in S] + \delta.
\]
A single $(\eps, \delta)$-pair is only a single point of the underlying
\emph{privacy profile} $\{(\eps,\delta(\eps)) : \eps \geq 0\}$~\cite{balle2020privacy,
kaissis2024beyond}, and mechanisms calibrated to the same point can have
substantially different profiles and so substantially different realised
adversarial risk. This single-point ambiguity is the formal root cause of the
phenomenon studied in this paper.

The \emph{Gaussian mechanism}, $\mathcal{M}(D) = f(D) + \mathcal{N}(0,\sigma^2 I)$
for a query $f$ of $L_2$ sensitivity $\Delta f$, satisfies $(\eps,\delta)$-DP
for all $(\eps, \delta)$ such that
$\Phi(\Delta f / 2\sigma - \eps\sigma/\Delta f) - e^\eps
\Phi(-\Delta f/2\sigma - \eps\sigma/\Delta f) \leq \delta$,
where $\Phi$ is the standard-normal CDF.
For training neural networks under DP, the standard procedure is DP-SGD with
Poisson subsampling~\cite{abadi2016deep}, in which each data point is included
in each minibatch independently with probability $p$, providing the well-known
\emph{subsampling amplification} of DP.

\subsection{Individual differential privacy and the sampling-based realisation}
\label{sec:idp}

\emph{Individual differential privacy}~\cite{jorgensen2015conservative}
relaxes the uniform $(\eps,\delta)$ guarantee to a per-record contract
$(\eps_i,\delta)$: the resulting algorithm must satisfy $(\eps_i,\delta)$-DP
\emph{with respect to record~$i$} for every $i$. Boenisch et
al.~\cite{boenisch2023have} translate per-record contracts into per-record
Poisson sampling probabilities while keeping a single noise multiplier $\sigma$
and a fixed expected batch size, which is the configuration we study.

\begin{definition}[Sampling-based iDP-SGD~\cite{boenisch2023have}]
\label{def:sampling-idp}
Let the dataset of size $M$ partition into $G$ \emph{privacy groups}
$\{G_1,\dots,G_G\}$ with sizes $|G_g|$. Each group has a target
$(\eps_g,\delta)$-DP budget. The mechanism trains by Poisson-subsampled
DP-SGD for $I$ iterations, where the noise multiplier $\sigma$ and the per-group
sampling rates $p_1,\dots,p_G$ are chosen such that
\begin{align}
\textnormal{(C1) for every } g, \quad
& \textnormal{the } I\textnormal{-fold } (p_g,\sigma)\textnormal{-subsampled} \notag\\
& \textnormal{Gaussian satisfies } (\eps_g,\delta)\textnormal{-DP at equality.} \notag\\
\textnormal{(C2)} \quad
& \sum_{g=1}^G \frac{|G_g|}{M}\,p_g \;=\; \frac{E_b}{M}, \label{eq:batch}
\end{align}
where $E_b$ is the desired expected batch size. We refer to $(\sigma,
p_1,\dots,p_G)$ as the \emph{mechanism solution} for the budget vector
$\boldsymbol{\eps} = (\eps_1,\dots,\eps_G)$.
\end{definition}

The single shared $\sigma$ in (C1) is the source of the privacy externality.
For any candidate $\sigma$, the largest sampling rate that keeps group $g$ at
its budget is a monotone-increasing function of $\sigma$ (more noise $\Rightarrow$
more room for sampling); the batch-size constraint (C2) then determines $\sigma$
as the unique value at which the per-group rates sum to $E_b$.  As a corollary,
the rate $p_g$ that group $g$ obtains is determined by the entire vector
$\boldsymbol{\eps}$, not by $\eps_g$ alone.

\subsection{MIA advantage and the privacy profile}

For a privacy-preserving mechanism $\mathcal{M}$ with trade-off function $T$
(equivalently, an $f$-DP function~\cite{dong2022gaussian}), the membership
inference advantage of the optimal Bayes-aware adversary
is~\cite{kaiser2026your}
\begin{equation}
\Adv(\mathcal{M}) \;=\; \sup_{\alpha \in [0,1]} \;\bigl(\,1 - \alpha - T(\alpha)\,\bigr),
\label{eq:adv}
\end{equation}
i.e.~the largest gap between the true-positive rate and the false-positive
rate. For the $I$-fold $(p, \sigma)$-subsampled Gaussian, $T$ can be
computed numerically by privacy-loss-distribution (PLD)
accounting~\cite{google-dp-accounting}, which is what we use throughout.
Crucially, $\Adv$ is determined by the \emph{full privacy profile}, not by any
single $(\eps,\delta)$ point.

\subsection{Game-theoretic terminology}

A finite normal-form game is a tuple
$\langle N, \Sigma, U\rangle$ with $N = \{1,\dots,n\}$ players, action sets
$\Sigma = \{S_1, \dots, S_n\}$, and payoff functions $U = \{u_1,\dots,u_n\}$.
A \emph{Nash equilibrium} (NE) is an action profile $\boldsymbol{s}^*$ such
that no player profits from a unilateral deviation:
$u_i(s_i^*,\boldsymbol{s}_{-i}^*) \geq u_i(s_i,\boldsymbol{s}_{-i}^*)$ for every
$s_i \in S_i$~\cite{nash1950equilibrium}. A \emph{best response} is the
utility-maximising action against a fixed strategy of the others. The
\emph{price of anarchy}~\cite{koutsoupias1999worst} compares the welfare of the
worst NE to the social optimum. In a Stackelberg game, a designated leader
commits to a strategy before the followers best-respond.
